\documentclass[aspectratio=169,10pt]{beamer}

\usetheme{metropolis}

\usepackage{amsmath,amssymb}
\usepackage{booktabs}
\usepackage{graphicx}
\usepackage{appendixnumberbeamer}

\title{Retrospective Climate Stressor\\Health Resilience Analysis}
\subtitle{Extending Burkart et al.\ (2021) to Measure\\Health System Adaptation to Observed Temperature Change}
\author{Aaron Osgood-Zimmerman\\{\small OZ Statistical Consulting}}
\institute{Prepared for Exemplars in Global Health, World Bank,\\Seed Global Health, and Management Sciences for Health}
\date{\today}

\begin{document}

\maketitle

% =====================================================================
\section{Background}
% =====================================================================

\begin{frame}{The Burkart et al.\ (2021) Framework}
  \begin{columns}[T]
    \begin{column}{\textwidth}
      \textbf{Two-part retrospective analysis:}
      \begin{enumerate}
        \item Estimate cause-specific exposure--response surfaces linking daily temperature to mortality (MR-BRT, 64.9M deaths, 9 countries)
        \item Apply globally to estimate attributable burden for 204 countries, 1990--2019
      \end{enumerate}

      \vspace{0.5em}
      \textbf{17 temperature-sensitive causes:}
      \begin{itemize}
       \item Selected after filtering out causes with no deaths and those with statistically insignificant relationships between RR and temperature
        \item Cardiovascular (IHD, stroke, hypertension, cardiomyopathy)
        \item Respiratory (LRI, COPD)
        \item Metabolic (diabetes, CKD)
        \item External (drowning, homicide, suicide, injuries, disasters)
      \end{itemize}
    \end{column}
  \end{columns}
\end{frame}



\begin{frame}{Key Concepts}
  \textbf{TMREL} --- Theoretical Minimum-Risk Exposure Level
  \begin{itemize}
    \item The daily temperature with the lowest death-weighted RR across all causes
    \item Location- and year-specific (varies with cause-of-death composition)
    \item RR curves are normalized so that $\text{RR} = 1.0$ at the TMREL
  \end{itemize}

  \vspace{0.5em}
  \textbf{PAF} --- Population Attributable Fraction, signifies the proportion of burden that could have potentially been prevented by removing/reducing the risk factor
  \begin{equation*}
    \text{PAF}_{cztd} = \frac{\text{RR}_{czt} - 1}{\text{RR}_{czt}}
  \end{equation*}
  {\small where $c$ = cause, $z$ = temperature zone, $t$ = daily temperature, $d$ = day, population-weighted across pixels $p$, summed across days $d$ within each year.} %Computed separately for heat (above TMREL) and cold (below TMREL).

  \vspace{0.5em}
  \textbf{Attributable burden} $B$ (Murray \& Lopez, 1999):
  \begin{equation*}
    B = M \times \text{PAF}
  \end{equation*}
  where $M$ = underlying cause-specific burden (deaths, YLLs, or DALYs)
\end{frame}

\begin{frame}{Exposure-Risk Function Calculations - Raw Data and Fit Curves}

  \centering
  \begin{columns}[T]
    \begin{column}{.6\textwidth}
      \includegraphics[width=\textwidth]{figures/fig-1a.png}
    \end{column}
    \begin{column}{.35\textwidth}
      \hspace{2in}
      {Figure 1a from Burkart et al.\ (2021). Raw data and Exposure--risk curves.}
      \hfill
      \end{column}
   \end{columns}
 \end{frame}

\begin{frame}{How the TMREL Is Calculated}
  \begin{columns}[T]
    \begin{column}{0.55\textwidth}
      The TMREL is the daily temperature that minimizes the \textbf{death-weighted average RR} across all 17 included causes:

      \vspace{0.3em}
      \begin{equation*}
        \overline{RR}_{ztly} = \frac{\displaystyle\sum_{c=1}^{17} \text{deaths}_{cztly} \cdot RR_{czt}}{\displaystyle\sum_{c=1}^{17} \text{deaths}_{cztly}}
      \end{equation*}

      \vspace{0.3em}
      \begin{equation*}
        \text{TMREL}_{zly} = \arg\min_t\; \overline{RR}_{ztly}
      \end{equation*}

      {\small where $c$ = cause, $z$ = temperature zone, $t$ = daily temperature, $l$ = location, $y$ = year.}
    \end{column}
    \begin{column}{0.42\textwidth}
      \textbf{Key properties:}
      \begin{itemize}
        \item Location- and year-specific
        \item Shifts when cause-of-death composition changes (e.g., if CVD declines relative to external causes)
        \item Computed from 1000 draws for uncertainty propagation
        \item Available for 1,034 GBD locations, 1990--2020
      \end{itemize}

      \vspace{0.3em}
      \textbf{For projections:} TMREL must be recomputed from projected cause-of-death composition
    \end{column}
  \end{columns}
\end{frame}

 \begin{frame}{Exposure-Risk Function Calculations - Rescaled}

  \centering
  \begin{columns}[T]
    \begin{column}{.55\textwidth}
      \includegraphics[width=\textwidth]{figures/fig-1b.png}
    \end{column}
    \begin{column}{.35\textwidth}
      \hspace{2in}
      {Figure 1b from Burkart et al.\ (2021). Exposure--risk curves rescaled so RR=1 at estimated TMREL.}
      \hfill
      \end{column}
   \end{columns}
\end{frame}


\begin{frame}{Heat, Cold, and Non-Optimal PAFs}

  Each pixel-day is classified relative to the TMREL:

  \vspace{0.3em}
  \begin{align*}
    \text{PAF}_{\text{heat}} &= \sum_{d:\, T_d > \text{TMREL}} \sum_p \frac{pop_p}{pop_l} \cdot \frac{\text{RR} - 1}{\text{RR}} \\[0.3em]
    \text{PAF}_{\text{cold}} &= \sum_{d:\, T_d < \text{TMREL}} \sum_p \frac{pop_p}{pop_l} \cdot \frac{\text{RR} - 1}{\text{RR}}
  \end{align*}
  {\small where $p$ = pixel, $l$ = location, $d$ = day. Days at exactly the TMREL are excluded ($\text{RR} = 1$).}

  \vspace{0.3em}
  \textbf{Non-optimal PAF} is the sum (heat and cold days are mutually exclusive):
  \begin{equation*}
    \text{PAF}_{\text{non-opt}} = \text{PAF}_{\text{heat}} + \text{PAF}_{\text{cold}}
  \end{equation*}

  \vspace{0.3em}
  \textbf{Key nuance:} Cold PAFs can be \alert{negative} for external causes (cold is protective against drowning, homicide). So $\text{PAF}_{\text{non-opt}}$ can be \emph{smaller} than $\text{PAF}_{\text{heat}}$ alone.
\end{frame}

\begin{frame}{Global Temperature-Attributable Burden (2019)}
  \begin{center}
    \includegraphics[width=0.85\textwidth]{figures/fig-4a.png}
  \end{center}
  \vspace{-0.5em}
  {\tiny Figure 4 from Burkart et al.\ (2021). Spatial distribution of all-cause DALYs (per 100,000) attributable to high temperature (A), low temperature (B), and non-optimal temperature (C) in 2019.}
\end{frame}

\begin{frame}{Global Temperature-Attributable Burden (2019)}
  \begin{center}
    \includegraphics[width=0.85\textwidth]{figures/fig-4b.png}
  \end{center}
  \vspace{-0.5em}
  {\tiny Figure 4 from Burkart et al.\ (2021). Spatial distribution of all-cause DALYs (per 100,000) attributable to high temperature (A) and low temperature (B) in 2019.}
\end{frame}

\begin{frame}{Global Temperature-Attributable Burden (2019)}
  \begin{center}
    \includegraphics[width=0.85\textwidth]{figures/fig-4c.png}
  \end{center}
  \vspace{-0.5em}
  {\tiny Figure 4 from Burkart et al.\ (2021). Spatial distribution of all-cause DALYs (per 100,000) attributable to high temperature (A) and low temperature (B) in 2019.}
\end{frame}




% \begin{frame}{Rising Heat Exposure Over Time}
%   \begin{center}
%     \includegraphics[width=0.7\textwidth]{figures/fig-3.png}
%   \end{center}
%   \vspace{-0.5em}
%   {\tiny Figure 3 from Burkart et al.\ (2021). Summary Exposure Values (SEVs) for high and low temperature, 1990--2019. Heat SEVs are steadily rising.}
% \end{frame}

% =====================================================================
\section{The Resilience Question}
% =====================================================================

\begin{frame}{From Burden Estimation to Resilience Measurement}

  \textbf{Burkart answers:} How large is the temperature-attributable health burden?

  \vspace{1em}
  \textbf{We ask:} How well are health systems \alert{absorbing} that burden as temperatures change?

  \vspace{1em}
  \textbf{Approach:} Compare two quantities over time:
  \begin{enumerate}
    \item The burden \textbf{expected} under observed rising temperatures, assuming population vulnerability remained at a historical baseline
    \item The burden \textbf{actually observed}
  \end{enumerate}

  \vspace{1em}
  The \alert{gap} between expected and observed burden is a measure of climate stressor health resilience --- applicable to any country with temperature data and cause-specific mortality.
\end{frame}

\begin{frame}{Attributable Burden = Burden $\times$ Exposure}

  The attributable burden $B$ in any year $t$ is:
  \begin{equation*}
    B(t) = \underbrace{M(t)}_{\text{underlying burden}} \times \underbrace{\text{PAF}\bigl(T(t)\bigr)}_{\text{exposure}}
  \end{equation*}
  where $M(t)$ = total cause-specific deaths (or YLLs/DALYs), $T(t)$ = daily temperature distribution in year $t$, and $\text{PAF}$ = population attributable fraction derived from the temperature exposure.

  \vspace{0.3em}
  Changes over time come from:
  \begin{itemize}
    \item \textbf{$M(t)$}: cause-specific mortality shifting due to health system improvements, epidemiological transition, population change
    \item \textbf{PAF$(T(t))$}: temperature distributions shifting due to climate change
  \end{itemize}

  \vspace{0.5em}
  To measure resilience, we \alert{disentangle} these two drivers.
\end{frame}

% =====================================================================
\section{Counterfactual Framework}
% =====================================================================

\begin{frame}{Four Scenarios}
  Fix a baseline period $t_0$ (e.g., first 2--3 years). Compute baseline mortality \emph{rates} $m_0$ (deaths per 100,000 by location $l$, cause $c$, age, sex). For each subsequent year $t$:

  \vspace{0.3em}
  \begin{table}
    \small
    \begin{tabular}{@{}llll@{}}
      \toprule
      \textbf{Scenario} & \textbf{Mortality} & \textbf{Temperature} & \textbf{Interpretation} \\
      \midrule
      $B_{\text{obs}}$ & Observed $M(t)$ & Observed $T(t)$ & What actually happened \\[0.3em]
      $B_{\text{noWarm}}$ & Observed $M(t)$ & Baseline $T_0$ & No warming counterfactual \\[0.3em]
      $B_{\text{noAdapt}}$ & $m_0 \times \text{Pop}(t)$ & Observed $T(t)$ & No adaptation counterfactual \\[0.3em]
      $B_{\text{base}}$ & $m_0 \times \text{Pop}(t)$ & Baseline $T_0$ & Demographic growth only \\
      \bottomrule
    \end{tabular}
  \end{table}

  \vspace{0.3em}
  {\small where $m_0$ = baseline mortality rates, $\text{Pop}(t)$ = population in year $t$, $T_0$ = baseline temperature distribution, $\text{PAF}_0$ = PAF computed from $T_0$.}
\end{frame}

\begin{frame}{Derived Quantities}

  \textbf{Excess burden from temperature change:}
  \begin{equation*}
    \Delta B_{\text{tempchange}}(t) = B_{\text{obs}}(t) - B_{\text{noWarm}}(t) = \Delta B_{\text{heat}} + \Delta B_{\text{cold}}
  \end{equation*}
  {\small Decomposes into heat (likely $+$) and cold (likely $-$ under warming) components.}

  \vspace{0.3em}
  \textbf{Burden absorbed by adaptation:}
  \begin{equation*}
    \Delta B_{\text{adapted}}(t) = B_{\text{noAdapt}}(t) - B_{\text{obs}}(t)
  \end{equation*}

  \vspace{0.5em}
  \textbf{Climate Stressor Health Resilience Index:}
  \begin{equation*}
    \boxed{\;R(t,l,c) = 1 - \frac{B_{\text{obs}}(t,l,c)}{B_{\text{noAdapt}}(t,l,c)}\;}
  \end{equation*}

  {\small where $l$ = location, $c$ = cause, $t$ = year.}
  \begin{itemize}
    \item $R = 0$: No adaptation (burden matches what baseline rates predict)
    \item $R > 0$: Positive adaptation (health system absorbing burden)
    \item $R < 0$: Maladaptation (burden exceeds prediction)
  \end{itemize}
\end{frame}

\begin{frame}{Decomposing the Change}
  \begin{columns}[T]
    \begin{column}{0.55\textwidth}
      Attributable burden is a product of three factors:
      \begin{equation*}
        B = \underbrace{m}_{\text{mortality rate}} \times \underbrace{\text{Pop}}_{\text{population}} \times \underbrace{\text{PAF}}_{\text{temperature}}
      \end{equation*}

      \vspace{0.2em}
      Das Gupta (1993) symmetric decomposition:
      \begin{align*}
        \Delta_m &= (m_1 - m_0) \cdot \frac{\text{Pop}_0 \!\cdot\! \text{PAF}_1 + \text{Pop}_1 \!\cdot\! \text{PAF}_0}{2} \\[0.2em]
        \Delta_{\text{Pop}} &= (\text{Pop}_1 - \text{Pop}_0) \cdot \frac{m_0 \!\cdot\! \text{PAF}_1 + m_1 \!\cdot\! \text{PAF}_0}{2} \\[0.2em]
        \Delta_{\text{PAF}} &= (\text{PAF}_1 - \text{PAF}_0) \cdot \frac{m_0 \!\cdot\! \text{Pop}_1 + m_1 \!\cdot\! \text{Pop}_0}{2}
      \end{align*}

      $\Delta_m + \Delta_{\text{Pop}} + \Delta_{\text{PAF}} = \Delta B$ \;\textit{exactly}
    \end{column}
    \begin{column}{0.42\textwidth}
      \textbf{Why decompose?}

      The resilience index $R$ tells us the health system is (or isn't) coping --- but not \emph{why}. Burden can change because:
      \begin{itemize}
        \item[$\Delta_m$:] Health systems improved (or worsened)
        \item[$\Delta_{\text{Pop}}$:] Population grew or aged
        \item[$\Delta_{\text{PAF}}$:] Temperatures shifted
      \end{itemize}

      \vspace{0.3em}
      \textbf{Policy value:} If $\Delta_{\text{PAF}} > 0$ but $|\Delta_m|$ is larger, health gains are outpacing warming. If not, the temperature signal is winning.
    \end{column}
  \end{columns}
\end{frame}

% =====================================================================
\section{Interpretation \& Strategy}
% =====================================================================

\begin{frame}{$\Delta B_{\text{tempchange}}$: Is Temperature Change a Problem Here?}
      \begin{equation*}
        \Delta B_{\text{tempchange}}(t) = B_{\text{obs}}(t) - B_{\text{noWarm}}(t)
        = M_{\text{obs}}(t) \times \bigl[\text{PAF}(T_{\text{obs}}) - \text{PAF}_0\bigr]
      \end{equation*}

      \vspace{0.3em}
      Measures the \alert{additional burden imposed by observed temperature change}, using actual mortality levels.

      \vspace{0.3em}
      \begin{itemize}
        \item $\Delta B_{\text{tempchange}} \approx 0$: Temperature shifts did not materially add burden
        \item $\Delta B_{\text{tempchange}} > 0$: Warming is imposing real additional burden
        \item $\Delta B_{\text{tempchange}} < 0$: Temperature change reduced burden (e.g., cold-climate locations benefiting from fewer cold days)
      \end{itemize}

      \vspace{0.5em}
      $\Delta B_{\text{tempchange}}$ answers: \emph{``Is temperature change adding burden here?''}\\
      $R$ answers: \emph{``Is the health system coping with it?''}
\end{frame}

\begin{frame}{Identifying Resilience Exemplars}

  Combining $\Delta B_{\text{tempchange}}$ with $R$ identifies three situations:

  \vspace{0.5em}
  \begin{itemize}
    \item \textbf{High $\Delta B_{tc}$, low $R$}: Stressor present, system \alert{not coping}\\
      {\small Priority for climate-health investment.}

    \vspace{0.5em}
    \item \textbf{High $\Delta B_{tc}$, high $R$}: \alert{Stressor Resilience Exemplar Candidate}\\
      {\small Temperature change is unambiguously adding burden, and the health system is unambiguously absorbing it. Cleanest exemplar signal --- study what these locations did differently.}

    \vspace{0.5em}
    \item \textbf{Large temp change, low $\Delta B_{tc}$}: \alert{Possible exemplar} --- needs verification\\
      {\small Could reflect extreme health system performance (underlying mortality dropped so far that even a real PAF increase produces little burden). Or could be geographic luck: location sits near the TMREL where heat and cold effects roughly cancel. Check the $\Delta B_{\text{heat}}$ / $\Delta B_{\text{cold}}$ split to distinguish.}
  \end{itemize}
\end{frame}

\begin{frame}{What Does ``Adaptation'' Capture?}

  The gap $B_{\text{noAdapt}} - B_{\text{obs}}$ reflects the net effect of changes in the underlying burden (for the 17 included causes), which may stem from:

  \begin{enumerate}
    \item \textbf{Temperature-specific adaptation}\\
      Early warning systems, cooling centers, AC, urban greening, behavioral change
    \item \textbf{General health system improvement}\\
      Better cardiovascular care, stroke treatment, diabetes management
    \item \textbf{Epidemiological transition}\\
      Shifting cause-of-death composition independent of climate
  \end{enumerate}

  \vspace{0.5em}
  \textbf{Partial identification strategies:}
  \begin{itemize}
    \item \textbf{Cross-cause:} $R$ high for IHD but low for drowning $\Rightarrow$ cardiovascular care, not climate-specific adaptation
    \item \textbf{Cross-location:} Similar temperature trends, different health capacity $\Rightarrow$ isolates health system contribution
  \end{itemize}
\end{frame}

% =====================================================================
\section{Data Requirements \& Implementation}
% =====================================================================

\begin{frame}{Data Requirements}

  \textbf{Available from Burkart/GBD (in hand):}
  \begin{itemize}
    \item Exposure--response curves (18 causes, 1000 draws each)
    \item TMRELs (1000 draws, 1034 locations, 1990--2020)
    \item GBD 2023 shapefile and location hierarchy
  \end{itemize}

  \vspace{0.5em}
  \textbf{Needed per country/location:}
  \begin{itemize}
    \item Daily temperature data (ERA5 for historical; SSP/CMIP6 for projections)
    \item Cause-specific mortality by age, sex, year (17 causes)
    \item Population estimates (gridded for pixel weighting)
    \item Life tables (for YLL/AVPP calculation)
  \end{itemize}

  \vspace{0.5em}
  \textbf{For projections to 2100} (companion project):
  \begin{itemize}
    \item Mortality projections by cause (from IHME, pending)
    \item Projected life tables (IHME or UN WPP)
    \item Projected TMRELs (computable from mortality projections)
  \end{itemize}
\end{frame}

\begin{frame}{Implementation Architecture}

  \textbf{Pipeline structure} (per location):
  \begin{enumerate}
    \item Extract/align temperature and population grids to GBD locations
    \item Assign temperature zones; compute pixel-day RRs from ERF curves
    \item Rescale RRs to 1.0 at year-specific TMREL
    \item Compute population-weighted PAFs (1000 draws)
    \item Multiply by cause-specific burden $\Rightarrow$ attributable deaths
    \item Compute counterfactuals and resilience index
  \end{enumerate}

  \vspace{0.5em}
  \textbf{Design:} Location-level parallelization with a wrapper script to run all locations. ERF curves and TMRELs are shared inputs; everything else is location-specific.
\end{frame}

% =====================================================================
\section{Policy Relevance}
% =====================================================================

\begin{frame}{Why This Matters}

  Shifts the question from \emph{``How bad is the burden?''} to\\
  \alert{``How well are health systems coping?''}

  \vspace{1em}
  \begin{itemize}
    \item \textbf{Identifies adaptation gaps:} locations and causes where $R$ is lowest are priority targets for climate-health investment
    \item \textbf{Quantifies adaptation value:} $\Delta B_{\text{adapted}}$ estimates what health improvements have already ``bought'' in climate resilience
    \item \textbf{Distinguishes general vs.\ climate-specific gains:} cross-cause analysis reveals whether to invest in health system capacity or targeted climate interventions
    \item \textbf{Uses only observed data:} no climate model outputs, emissions scenarios, or assumptions about future adaptation required
  \end{itemize}

  \vspace{1em}
  The resilience analysis also serves as a \textbf{validation step} for the companion burden projection work.
\end{frame}

% =====================================================================
\section{Extension to Other Stressors}
% =====================================================================

\begin{frame}{Framework Requirements for Other Stressors}

  The counterfactual structure extends to \alert{slow-onset, continuous stressors} that satisfy:

  \vspace{0.5em}
  \begin{enumerate}
    \item \textbf{Continuous exposure distribution} that shifts gradually over time\\
      {\small (e.g., daily PM\textsubscript{2.5} concentrations, not binary flood occurrence)}
    \item \textbf{Smooth dose--response relationship} mapping exposure to cause-specific RR\\
      {\small (enables PAF calculation as an integral over the distribution)}
    \item \textbf{Gradual temporal change} driven by climate trend, not stochastic events
  \end{enumerate}

  \vspace{0.5em}
  \textbf{Poor fit:} Acute climate shocks (floods, storms, droughts) --- binary exposure, event-level dose--response, stochastic inter-annual variation. These need different methods (event attribution, impact evaluation).
\end{frame}

\begin{frame}{Candidate Stressors}
  \textbf{Strong fit} --- GBD infrastructure exists:
  \begin{itemize}
    \item \textbf{PM\textsubscript{2.5}}: Wildfires, stagnation; GBD integrated exposure--response curves
    \item \textbf{Ground-level ozone}: Temperature accelerates formation; GBD estimates ozone$\to$COPD
    \item \textbf{UV radiation}: Ozone/cloud changes; established dose--response for skin cancer, cataracts
  \end{itemize}

  \vspace{0.3em}
  \textbf{Moderate fit} --- exposure--response less mature:
  \begin{itemize}
    \item \textbf{Ambient humidity}: Independent CVD/respiratory risk factor; hard to disentangle from temperature
    \item \textbf{Pollen/aeroallergens}: Longer seasons, higher concentrations; asthma outcomes (morbidity, not mortality)
    \item \textbf{Vector-borne disease range}: Continuous transmission suitability, but heavily mediated by non-climate factors (bed nets, treatment)
  \end{itemize}

  \vspace{0.3em}
  {\small Note: For stressors with primarily non-fatal outcomes (UV, pollen), use YLDs or DALYs as $M$ rather than deaths --- the $B = M \times \text{PAF}$ framework accommodates any GBD burden measure.}
\end{frame}

\begin{frame}{Toward a Composite Climate Stressor Health Resilience Index}

  \textbf{Key insight:} The same cause-specific mortality data ($M$) can serve as the underlying burden for multiple stressors simultaneously
  \begin{itemize}
    \item GBD already does this: temperature and PM\textsubscript{2.5} both produce PAFs for IHD, applied to the same death counts
    \item Each stressor needs its own exposure data and exposure--response curves
    \item The resilience index $R(t,l,c)$ and decomposition apply without modification
  \end{itemize}

  \vspace{0.5em}
  \textbf{Composite index:} Combine stressor-specific resilience indices, weighted by relative burden contribution:
  \begin{itemize}
    \item Single summary: how well does a health system cope with the aggregate slow-onset climate threat?
    \item Stressor-specific components: where is targeted investment most needed?
  \end{itemize}
\end{frame}

\begin{frame}[standout]
  Questions?
\end{frame}

% =====================================================================
\appendix
% =====================================================================

\begin{frame}{References}
  \small
  \begin{enumerate}
    \item Burkart KG, Brauer M, Aravkin AY, et al.\ Estimating the cause-specific relative risks of non-optimal temperature on daily mortality: a two-part modelling approach applied to the Global Burden of Disease Study.\ \textit{Lancet}. 2021;398:685--697.
    \item Murray CJL, Lopez AD.\ On the comparable quantification of health risks: lessons from the Global Burden of Disease Study.\ \textit{Epidemiology}. 1999;10:594--605.
    \item Das Gupta P.\ Standardization and decomposition of rates: a user's manual.\ \textit{U.S.\ Bureau of the Census, Current Population Reports}. 1993;P23--186.
    \item Kitagawa EM.\ Components of a difference between two rates.\ \textit{J Am Stat Assoc}. 1955;50(272):1168--1194.
  \end{enumerate}
\end{frame}

\begin{frame}{PAF and Burden Estimates (Burkart Table 1)}
  \begin{center}
    \includegraphics[width=0.85\textwidth]{figures/table1-06.png}
  \end{center}
  \vspace{-0.5em}
  {\tiny Table from Burkart et al.\ (2021). PAF and burden estimates for high, low, and non-optimal temperature by country, 1990 and 2019.}
\end{frame}

\begin{frame}{Global Non-Optimal Temperature Burden (2019)}
  \begin{center}
    \includegraphics[width=0.85\textwidth]{figures/fig4c-10.png}
  \end{center}
  \vspace{-0.5em}
  {\tiny Figure 4C from Burkart et al.\ (2021). Spatial distribution of all-cause DALYs (per 100,000) attributable to non-optimal temperature in 2019.}
\end{frame}

\end{document}
